\documentclass[12pt]{article}
\usepackage[utf8]{inputenc}
\usepackage[T1]{fontenc}
\usepackage[margin=1in]{geometry}
\usepackage{setspace}
\usepackage{xcolor}
\usepackage[colorlinks=true, urlcolor=blue!60!black, linkcolor=blue!60!black]{hyperref}
\setlength{\parindent}{2em}
\setlength{\parskip}{1.25em}
\renewcommand{\baselinestretch}{1.0}
\title{\textit{SEC v. Jarkesy}, 603 U.S. 109 (2024)}
\author{Case Brief}
\date{}

\begin{document}
\maketitle

\section{Case Metadata}
\begin{description}
  \item[Citation:] 603~U.S.~109 (2024) (No.~22--859).
  \item[Argued / Decided:] Argued November~29, 2023; Decided June~27, 2024.
  \item[Vote:] 6--3 (affirming and remanding 34~F.4th~446 (5th~Cir.~2022)).
  \item[Author:] Chief Justice Roberts (joined by Justices Thomas, Alito, Gorsuch, Kavanaugh, and Barrett).
  \item[Concurrences / Dissents:] Justice Gorsuch filed a concurring opinion (joined by Justice Thomas); Justice Sotomayor filed a dissenting opinion (joined by Justices Kagan and Jackson).
  \item[Chapter / Topic:] Separation of Powers and the Executive --- Article~III Judicial Power; Seventh Amendment Right to Civil Jury Trial; Public Rights Exception; Administrative Adjudication.
  \item[Primary Source URLs:] \href{https://www.supremecourt.gov/opinions/23pdf/603us1r50_7kh7.pdf}{SupremeCourt.gov Page Proof} (\href{https://www.supremecourt.gov/opinions/23pdf/22-859_1924.pdf}{Slip Opinion}); \href{https://supreme.justia.com/cases/federal/us/603/22-859/}{Justia U.S. Supreme Court Center}.
\end{description}

\section{Facts and Procedural History}
Between 2007 and 2010, George~R.~Jarkesy, Jr. launched two investment funds that raised approximately \$24~million from 120 accredited investors, with Patriot28, LLC serving as the funds' investment adviser. Prior to 2010, the Securities and Exchange Commission (SEC) could seek civil monetary penalties against unregistered entities only by filing suit in federal district court. In \S~929P(a) of the Dodd-Frank Wall Street Reform and Consumer Protection Act of 2010, 124~Stat.~1862--1864, Congress expanded the SEC's authority so that it could seek civil penalties of up to \$725,000 per violation either in federal court or through its own in-house administrative proceedings. In federal court, an Article~III judge presides, a jury finds the facts, and the Federal Rules of Evidence and ordinary discovery rules apply. In SEC in-house proceedings, by contrast, there is no jury, the Commission or an Administrative Law Judge (ALJ) finds facts and resolves discovery disputes under the SEC's Rules of Practice (which permit hearsay), and subsequent judicial review treats agency factual findings as ``conclusive'' if supported by substantial evidence.

In 2013, relying on its Dodd-Frank authority, the SEC initiated an in-house administrative enforcement action against Jarkesy and Patriot28 under the antifraud provisions of the Securities Act of 1933 (\S~17(a)), the Securities Exchange Act of 1934 (\S~10(b) and Rule~10b-5), and the Investment Advisers Act of 1940 (\S~206 and Rule~206(4)-8). The SEC alleged that respondents misled investors by misrepresenting their investment strategies, lying about the identity of the funds' auditor and prime broker, and inflating the funds' claimed value to collect larger management fees. An SEC ALJ issued an initial decision finding liability in 2014. After six years of administrative review, the full Commission issued a final order in 2020 imposing a \$300,000 civil penalty, a cease-and-desist order, an industry and penny-stock bar against Jarkesy, and disgorgement of \$685,000 against Patriot28.

Jarkesy and Patriot28 petitioned the Fifth Circuit for review. A divided panel vacated the SEC's order on three independent constitutional grounds, 34~F.4th~446 (5th~Cir.~2022) (Judge Davis dissenting): (1) adjudicating civil-penalty securities-fraud claims in-house violated the Seventh Amendment right to a jury trial; (2) Congress violated the nondelegation doctrine by giving the SEC unguided discretion to choose between federal court and administrative adjudication; and (3) two layers of for-cause removal protection insulating SEC ALJs violated Article~II. After the Fifth Circuit denied rehearing en banc, 51~F.4th~644 (2022), the Supreme Court granted certiorari on June~30, 2023.

\section{Issue(s) Presented}
Whether the Seventh Amendment entitles a defendant to a jury trial in an Article~III court when the SEC seeks civil monetary penalties for alleged violations of the federal securities laws' antifraud provisions. (Because the Court resolved the case on the Seventh Amendment ground alone, it expressly declined to reach the Fifth Circuit's nondelegation and Article~II removal holdings.)

\section{Holding and Rule}
\textbf{Holding:} Yes. When the SEC seeks civil penalties against a defendant for securities fraud, the Seventh Amendment entitles the defendant to a jury trial in an Article~III court.

\textbf{Rule:} Under the two-step framework of \textit{Tull v. United States}, 481~U.S.~412 (1987), and \textit{Granfinanciera, S.A. v. Nordberg}, 492~U.S.~33 (1989), a statutory enforcement action implicates the Seventh Amendment if the claim is ``legal in nature''---with the punitive, non-compensatory nature of civil monetary penalties being ``all but dispositive''---and Congress may not assign such a claim to a non-Article~III agency tribunal unless it falls within one of the narrow, historically grounded categories of the ``public rights'' exception.

\section{Reasoning}
Writing for a six-Justice majority, Chief Justice Roberts applied a two-step constitutional inquiry:
\begin{enumerate}
  \item \textbf{Step One --- Implication of the Seventh Amendment:} The Seventh Amendment preserves the jury-trial right in all statutory claims that are ``legal in nature,'' evaluated by both the cause of action and---more importantly---the remedy sought.
  \begin{itemize}
    \item \textit{Remedy (All But Dispositive):} Monetary relief is legal rather than equitable when designed to punish or deter rather than solely to ``restore the status quo.'' The securities statutes condition civil penalties on six factors (including fraud, culpability, recidivism, and deterrence) and three escalating tiers tied to scienter rather than victim loss, and the SEC is not required to distribute penalty proceeds to injured investors. Accordingly, SEC civil penalties are a prototypical common-law legal remedy enforceable only in courts of law.
    \item \textit{Cause of Action:} Federal securities fraud shares a close ancestral relationship with common-law fraud: both target the misrepresentation or concealment of material facts, and Congress deliberately transplanted common-law terms of art (``fraud,'' ``deceit,'' ``manipulative'') that bring their ``old soil'' with them.
  \end{itemize}
  \item \textbf{Step Two --- Inapplicability of the Public Rights Exception:} Under \textit{Murray's Lessee v. Hoboken Land \& Improvement Co.}, 18~How.~272 (1856), and \textit{Stern v. Marshall}, 564~U.S.~462 (2011), Congress may not ``withdraw from judicial cognizance any matter which, from its nature, is the subject of a suit at the common law.'' Although the Court has recognized an unwritten ``public rights'' exception for historic categories of sovereign prerogatives (revenue collection, customs enforcement, immigration, relations with Indian tribes, administration of public lands, and the granting of public benefits), the presumption runs strongly in favor of Article~III courts.
  \item \textbf{Application of \textit{Granfinanciera} and Narrowing of \textit{Atlas Roofing}:} Under \textit{Granfinanciera}, Congress cannot ``conjure away the Seventh Amendment by mandating that traditional legal claims be \ldots\ taken to an administrative tribunal.'' That the federal government is the prosecuting party does not itself trigger the public rights exception; ``what matters is the substance of the suit, not where it is brought, who brings it, or how it is labeled.'' The majority distinguished \textit{Atlas Roofing Co. v. OSHRC}, 430~U.S.~442 (1977), characterizing the Occupational Safety and Health Act's technical workplace standards as a self-consciously novel regulatory code ``unknown to the common law'' rather than a claim rooted in common-law fraud, while noting that \textit{Tull} and \textit{Granfinanciera} foreclose reading \textit{Atlas Roofing} as a blanket license for agency civil-penalty adjudication.
\end{enumerate}

\section{Concurring and Dissenting Opinions}
\begin{itemize}
  \item \textbf{Justice Gorsuch, concurring (joined by Justice Thomas):} Emphasized that the Seventh Amendment operates in tandem with Article~III's guarantee of an independent judge and the Fifth Amendment's Due Process Clause to secure a ``fair trial in a fair tribunal.'' He drew a direct historical parallel between modern SEC in-house proceedings (where the SEC reportedly won roughly 90\% of contested cases versus 69\% in federal court) and colonial British vice-admiralty courts, urged anchoring the public rights exception strictly to historically established categories with unbroken Founding-era pedigrees, and criticized \textit{Atlas Roofing} as an atextual outlier.
  \item \textbf{Justice Sotomayor, dissenting (joined by Justices Kagan and Jackson):} Condemned the decision as a ``power grab'' that upends nearly two centuries of precedent (\textit{Murray's Lessee}, \textit{Oceanic Steam Navigation Co. v. Stranahan}, \textit{Crowell v. Benson}, and a unanimous \textit{Atlas Roofing}). She argued that the threshold question is Article~III, not the Seventh Amendment, and that under settled doctrine a statutory civil-penalty claim belonging to the federal government in its sovereign capacity is the ``heartland'' of public rights. She faulted the majority for misapplying \textit{Tull} (a federal-court case) and \textit{Granfinanciera} (a private-party bankruptcy dispute) and warned that the ruling imperils over 200 statutes across more than two dozen federal agencies (including OSHA, FERC, USDA, EPA, FCC, and CFTC), many of which lack statutory authority to sue in federal court.
\end{itemize}

\section{Key Quotations}
\begin{quote}
``A defendant facing a fraud suit has the right to be tried by a jury of his peers before a neutral adjudicator. Rather than recognize that right, the dissent would permit Congress to concentrate the roles of prosecutor, judge, and jury in the hands of the Executive Branch. That is the very opposite of the separation of powers that the Constitution demands.'' \\
--- Chief Justice Roberts, \href{https://supreme.justia.com/cases/federal/us/603/22-859/}{603~U.S.~at~140--141} (Majority Opinion).
\end{quote}

\begin{quote}
``In this case, the remedy is all but dispositive. \ldots\ In sum, the civil penalties in this case are designed to punish and deter, not to compensate. They are therefore `a type of remedy at common law that could only be enforced in courts of law.' '' \\
--- Chief Justice Roberts, \href{https://supreme.justia.com/cases/federal/us/603/22-859/}{603~U.S.~at~123, 125} (Majority Opinion) (quoting \textit{Tull v. United States}, 481~U.S.~at~422).
\end{quote}

\begin{quote}
``What matters is the substance of the suit, not where it is brought, who brings it, or how it is labeled.'' \\
--- Chief Justice Roberts, \href{https://supreme.justia.com/cases/federal/us/603/22-859/}{603~U.S.~at~135} (Majority Opinion).
\end{quote}

\begin{quote}
``The Seventh Amendment's jury-trial right does not work alone. It operates together with Article~III and the Due Process Clause of the Fifth Amendment to limit how the government may go about depriving an individual of life, liberty, or property.'' \\
--- Justice Gorsuch, \href{https://supreme.justia.com/cases/federal/us/603/22-859/}{603~U.S.~at~141} (Concurring Opinion).
\end{quote}

\begin{quote}
``Today, for the very first time, this Court holds that Congress violated the Constitution by authorizing a federal agency to adjudicate a statutory right that inheres in the Government in its sovereign capacity, also known as a public right. \ldots\ Make no mistake: Today's decision is a power grab.'' \\
--- Justice Sotomayor, \href{https://supreme.justia.com/cases/federal/us/603/22-859/}{603~U.S.~at~167, 202} (Dissenting Opinion).
\end{quote}

\section{Doctrinal Significance and Modern Legacy (Through 2026)}
Decided one day before \href{https://supreme.justia.com/cases/federal/us/603/22-451/}{\textit{Loper Bright Enterprises v. Raimondo}, 603~U.S.~369 (2024)} (which overruled \textit{Chevron} deference), \textit{Jarkesy} marked a landmark reassertion of Article~III judicial power and the Seventh Amendment against executive-branch adjudication. Immediately after \textit{Jarkesy}, regulated parties launched Seventh Amendment challenges against administrative penalty proceedings across the federal government---including the FCC, FERC, NLRB, DOL, CFTC, and USDA (with a 2025 study finding that the USDA sharply curtailed civil penalties under the Animal Welfare Act in favor of nonbinding warnings; see \href{https://en.wikipedia.org/wiki/SEC_v._Jarkesy}{post-\textit{Jarkesy} enforcement overview}).

Two major developments through 2026 have defined \textit{Jarkesy}'s practical reach:
\begin{itemize}
  \item \textbf{Clarifying the Outer Boundary --- \textit{FCC v. AT\&T, Inc.} (2026):} After the Fifth Circuit applied \textit{Jarkesy} to vacate a \$57~million FCC forfeiture order against AT\&T (\textit{AT\&T v. FCC}, 149~F.4th~491 (5th~Cir.~2025)), creating a split with the Second Circuit (\textit{Verizon Communications, Inc. v. FCC}, 156~F.4th~86 (2d~Cir.~2025)) and D.C.~Circuit (\textit{Sprint Corp. v. FCC}, 151~F.4th~347 (D.C.~Cir.~2025)), the Supreme Court resolved the conflict in \href{https://supreme.justia.com/cases/federal/us/608/25-406/}{\textit{FCC v. AT\&T, Inc.}, 608~U.S.~\_\_\_ (June~4, 2026)}. In an 8--1 opinion by Chief Justice Roberts (with only Justice Thomas dissenting), the Court reversed the Fifth Circuit and upheld the FCC's forfeiture framework under 47~U.S.C.~\S~503(b)(4). Distinguishing \textit{Jarkesy}---where SEC penalty orders were immediately enforceable via garnishment or tax offset without any jury review---the Court held that FCC forfeiture orders create no binding obligation to pay and carry no preclusive factual weight; to collect, DOJ must file a \S~504(a) enforcement suit in federal district court where the carrier receives a full \textit{trial de novo} before a jury (``Before a regulated party can be made to pay, the jury gets the last word''). Lower courts had already reached similar conclusions for FERC penalty schemes affording \textit{de novo} district court review (\textit{American Efficient LLC v. FERC}, M.D.N.C.~Nov.~24, 2025; see \href{https://www.eba-net.org/wp-content/uploads/2026/05/EBA-Brief-2026-Vol-7.1.pdf}{Energy Bar Ass'n Brief, May~2026}).
  \item \textbf{Unresolved Article~II Removal Issues:} Because \textit{Jarkesy} declined to reach the Fifth Circuit's ALJ removal holding, Article~II removal challenges migrated to separate litigation in 2025--2026---most notably \textit{Space Exploration Technologies Corp. v. NLRB} (5th~Cir.~2025) and Supreme Court proceedings addressing the continued vitality of \textit{Humphrey's Executor} in \textit{Trump v. Slaughter} and \textit{Trump v. Cook} (\href{https://www.eba-net.org/wp-content/uploads/2026/05/EBA-Brief-2026-Vol-7.1.pdf}{Energy Bar Ass'n Brief, May~2026}).
\end{itemize}

\section{Pedagogical Critiques and Exam Hypotheticals}
\textbf{Critiques:} Scholars and the dissent highlight two doctrinal tensions in \textit{Jarkesy}. First, the majority's statement that ``the remedy is all but dispositive'' sits uneasily beside its effort to preserve \textit{Atlas Roofing}, because OSHA civil penalties were equally punitive and non-compensatory; if punitive civil penalties trigger the Seventh Amendment unless a historical public-rights exception applies, \textit{Atlas Roofing} survives on razor-thin footing. Second, by refusing to articulate a comprehensive test for ``public rights,'' the Court left lower courts struggling to classify modern regulatory wrongs (such as environmental, labor, or data-privacy violations) that bear only loose analogies to common-law torts like nuisance or negligence. Conversely, proponents praise \textit{Jarkesy} (especially as refined by \textit{FCC v. AT\&T}) for restoring a clear constitutional baseline: agencies may investigate and issue preliminary assessments, but the Executive cannot unilaterally levy binding punitive fines for common-law-like wrongs without affording an Article~III jury trial.

\textbf{Exam Hypothetical:} Congress amends the Clean Air Act to empower the Environmental Protection Agency (EPA) to impose up to \$500,000 per violation in immediately enforceable civil penalties---collectible via Treasury offset without a federal-court collection suit---against factories that emit pollutants exceeding complex numerical parts-per-million thresholds. An ALJ adjudicates liability, and judicial review is limited to substantial-evidence review in the court of appeals. A penalized factory challenges the order under the Seventh Amendment.
\begin{itemize}
  \item \textbf{Challenger's Argument (Under \textit{Jarkesy}):} Because the EPA's civil penalties are punitive, non-compensatory, and immediately enforceable without a \textit{de novo} jury trial (distinguishing \textit{FCC v. AT\&T}), the remedy is ``all but dispositive.'' Moreover, air-pollution enforcement is closely analogous to common-law public and private nuisance, and environmental regulation under the Commerce Clause is not one of \textit{Jarkesy}'s historic public-rights categories (revenue, customs, immigration, public lands, Indian affairs, or public benefits).
  \item \textbf{Opposite Side Contrast (Distinguishing \textit{Jarkesy} via \textit{Atlas Roofing}):} The EPA will argue that, unlike the SEC's antifraud provisions (which explicitly borrowed common-law ``fraud'' terms of art), numerical parts-per-million emission thresholds resemble the technical trench-sloping ``building code'' standards upheld in \textit{Atlas Roofing}---prophylactic statutory obligations ``unknown to the common law'' that require no proof of private injury or nuisance. Furthermore, if Congress instead structured the EPA order as nonbinding unless enforced by DOJ in a \textit{de novo} federal district court jury trial, the scheme would be squarely constitutional under \textit{FCC v. AT\&T} (2026).
\end{itemize}

\end{document}
